% Part: computability
% Chapter: computability-theory
% Section: ce-sets

\documentclass[../../../include/open-logic-section]{subfiles}

\begin{document}

\olfileid{cmp}{thy}{ces}
\olsection{په محاسبوي ډول د شمېر وړ سټونه}

\begin{defn}
يو سټ هغه وخت \emph{په محاسبوي ډول د شمېر وړ} دے چې تش وي يا د يوې
محاسبه کېدونکې تابع د قيمتونو سټ وي۔
\end{defn}

\begin{history}
په محاسبوي ډول د شمېر وړ سټونه د دې پر ځاے \emph{په بازګشتي ډول د
  شمېر وړ} سټونه هم بلل کېږي۔ دا اصلي اصطلاح ده، او نن دواړه نومونه،
او ورسره لنډونونه ``c.e.'' او ``r.e.''، په عام ډول کارېږي۔
\end{history}

\begin{explain}
بايد په دې فکر وکړئ چې تعريف څه مانا لري او اصطلاح ولې مناسبه ده۔ نظر
دا دے چې که~$S$ د محاسبه کېدونکې تابع~$f$ د قيمتونو سټ وي، نو
\[
S = \{ f(0), f(1), f(2), \dots \},
\]
او ځکه~$f$ د~$S$ د غړو د ``شمېرلو'' په توګه ګڼل کېدے شي۔ پام وکړئ
چې د تعريف له مخې~$f$ لازمه نۀ ده زياتېدونکې تابع وي، يعنې شمېرنه
لازمه نۀ ده په زياتېدونکي ترتيب وي۔ په حقيقت کښې~$f$ لازمه نۀ ده
يو په يو هم وي، يعنې د~$S$ په شمېرنه $f(0)$، $f(1)$، $f(2)$،
\dots{} کښې تکرار منل شوے دے۔ د بېلګې په توګه، ثابته تابع
$f(x) = 0$ سټ~$\{ 0 \}$ شمېري۔
\end{explain}

هر محاسبه کېدونکے سټ په محاسبوي ډول د شمېر وړ دے۔ د دې د ليدلو
دپاره فرض کړئ~$S$ محاسبه کېدونکے دے۔ که~$S$ تش وي، نو د تعريف له
مخې په محاسبوي ډول د شمېر وړ دے۔ که نه، $a$ د~$S$ هر يو غړے وي۔
تابع~$f$ داسې تعريف کړئ:
\[
f(x) =
\begin{cases}
x & \text{که $\Char{S}(x) = 1$ وي} \\
a & \text{که داسې نۀ وي۔}
\end{cases}
\]
بيا~$f$ محاسبه کېدونکې تابع ده، او~$S$ د~$f$ د قيمتونو سټ دے۔

\end{document}
